\documentclass[11pt,letterpaper]{article}
\usepackage[spanish]{babel}
\usepackage{fontspec}
\usepackage[margin=2.1cm]{geometry}
\usepackage{graphicx,booktabs,longtable,array,xcolor,hyperref,fancyhdr,enumitem}
\setmainfont{Georgia}
\setsansfont{Arial}
\definecolor{azul}{HTML}{173B59}
\definecolor{acento}{HTML}{14A6B5}
\hypersetup{colorlinks=true,linkcolor=azul,urlcolor=azul}
\setlength{\parskip}{5pt}
\setlength{\parindent}{0pt}
\pagestyle{fancy}\fancyhf{}\fancyhead[L]{\small CGICH \textbar\ Manual de usuario}\fancyhead[R]{\small 2027-1}\fancyfoot[C]{\thepage}
\newcommand{\captura}[3]{\begin{figure}[htbp]\centering\includegraphics[width=#2\linewidth,height=0.32\textheight,keepaspectratio]{../capturas/#1}\caption{#3}\end{figure}}
\begin{document}
\begin{titlepage}
\color{azul}\sffamily
\vspace*{2.8cm}
{\Large UNIVERSIDAD NACIONAL AUTÓNOMA DE MÉXICO\par}
{\large Facultad de Ingeniería \textbar\ Computación Gráfica e Interacción Humano-Computadora\par}
\vspace{2cm}
{\Huge\bfseries Manual de usuario\par}
\vspace{0.5cm}{\Large Viajes espaciales y evolución de especies intergalácticas\par}
\vspace{1.5cm}\rule{\linewidth}{2pt}\vspace{0.5cm}
{\large David Díaz Antúnez\par Jonhatan Alexis Martínez Ojeda\par Venancio Miguel Tapia Verástegui\par}
\vfill{\large Semestre 2027-1 \textbar\ 7 de noviembre de 2026\par}
{\small Documentado a partir del candidato Release v0.4.0-rc.2.\par}
\end{titlepage}
\tableofcontents\newpage

\section{Antes de comenzar}
Esta experiencia contiene una base espacial, una pista de carrera, el planeta rocoso del Vacío, el planeta selvático Shinobi y una zona oceánica. Se puede recorrer a pie con Rek'Sai, competir en una nave, volar entre planetas y observar eventos desde una cámara orbital. El indicador de la esquina superior muestra el modo, la zona y el objeto activo.

\subsection{Instalación y arranque}
Se necesita Windows x64 y una tarjeta o controlador compatible con OpenGL 3.3. Descargue el ZIP completo de Release, compruebe su hash SHA-256 si se entregó uno y extráigalo conservando su estructura. Entre a la carpeta extraída y ejecute \texttt{Jugar.cmd}. No mueva solo el archivo EXE: el juego necesita sus carpetas \texttt{Models/}, \texttt{Textures/}, \texttt{shaders/} y las bibliotecas incluidas. \texttt{Verify-Release.ps1} verifica los archivos del paquete; \texttt{Jugar.cmd --probar-release} recorre una prueba automática. Para jugar, use \texttt{Jugar.cmd} sin argumentos.

Al abrirse, el personaje comienza en la base espacial. Haga clic dentro de la ventana si las teclas no responden. \textbf{Esc} cierra el programa.

\section{Controles}
\small
\begin{longtable}{p{0.20\linewidth}p{0.72\linewidth}}
\toprule Tecla o dispositivo & Acción \\ \midrule\endhead
W, A, S, D & Mover a Rek'Sai o la nave. En carrera, W/S aceleran y frenan; A/D cambian de carril.\\
Ratón & Orientar la vista en planeta, carrera y vuelo.\\
F1 & Recorrido de planeta en tercera persona. Desde espacio vuelve al último planeta visitado.\\
F2 & Recorrido de carrera: toma la nave del jugador y sigue la pista.\\
F3 & Vuelo en tercera persona entre planetas. W/A/S/D desplazan la nave; Q desciende y E asciende en espacio.\\
F4 & Observación orbital de lanzamiento, hangar, planetas y satélite.\\
Tab / rueda & En observación, cambiar de objetivo / acercar o alejar el encuadre.\\
F5 / F6 / F7 / F8 & Accesos a selva / roca / océano / espacio.\\
J & Abordar o desembarcar junto a la nave.\\
E & Interactuar cerca de un planeta, o regresar desde su superficie; en vuelo espacial, mantener para ascender.\\
1 & Iniciar el lanzamiento del cohete cuando la ubicación o el objetivo observado lo permiten. Al estar en órbita, reinicia su secuencia.\\
2 / 3 & Abrir o cerrar la puerta de la cápsula / paneles del satélite.\\
4 / 5 & Lanzar Fighter\_02 desde el hangar / abrir o cerrar la puerta del hangar.\\
6 / 7 & En Rek'Sai: reproducir el cohete de Jona / reiniciarlo.\\
F / L / K & Alternar la linterna / faro de nave / foco de base.\\
C / F9 & Alternar cámara libre / liberar o capturar el cursor.\\
Esc & Salir.\\\bottomrule
\end{longtable}\normalsize
Las teclas de acción actúan según la ubicación y el modo. Si no ocurre nada, revise el texto del indicador y pruebe desde el objetivo correspondiente.

\section{Recorridos guiados}
\subsection{R1: caminar por los planetas}
Pulse \textbf{F5} para llegar a la selva o \textbf{F6} para llegar al planeta rocoso; luego \textbf{F1}. Si aparece en la nave, use \textbf{J} para bajar. Camine con WASD y gire la vista con el ratón. La cámara sigue a Rek'Sai sobre el plano del terreno. Observe criaturas, árboles, rocas, cristales, portales y las diferencias de ambiente entre ambos mundos. La zona oceánica se visita con \textbf{F7}. \textbf{E} permite volver al espacio desde las superficies.
\captura{fauna_vacio.png}{0.83}{Criaturas y ambientación del planeta rocoso. La imagen reúne cuatro vistas de su fauna.}

\subsection{R2: carrera espacial}
Pulse \textbf{F2}. La cámara se coloca detrás de la nave del jugador. Mantenga \textbf{W} para avanzar, \textbf{S} para reducir velocidad y \textbf{A/D} para cambiar de carril. Otras dos naves recorren continuamente el circuito, que enlaza espacio y tramos de planetas. Si parte a pie en la base, también puede abordar la nave cercana con \textbf{J}.
\captura{r2_carrera_2026-11-06.png}{0.87}{Nave del jugador sobre el circuito en el modo de carrera.}

\subsection{R3: vuelo entre planetas}
Pulse \textbf{F3}. El ratón orienta la mirada y la nave; WASD desplaza, \textbf{Q} baja y \textbf{E} sube en espacio. Aproxímese a los cuerpos celestes. Puede entrar a un planeta al cruzar su atmósfera o utilizar los accesos F5/F6. \textbf{F8} vuelve a la zona espacial.

\subsection{R4: observar lanzamiento y estación}
Pulse \textbf{F4}. La primera vista disponible enfoca el cohete de Shinobi. Pulse \textbf{1} para iniciar ignición y ascenso. Con \textbf{Tab} pase al hangar, los planetas, el satélite y el cohete de la base. La rueda cambia la distancia de observación. El personaje conserva su ubicación mientras la cámara muestra otros lugares. F1, F2 o F3 reanudan un recorrido interactivo.
\captura{cohete_20261106_shinobi_ascenso.png}{0.83}{Lanzamiento del cohete en Shinobi; el estado aparece en el indicador.}
\captura{r4_hangar_2026-11-06.png}{0.83}{Estación espacial y Fighter\_02 en observación.}

\section{Objetos interactivos y ambiente}
La tecla \textbf{2} alterna la puerta de la cápsula y \textbf{3} los paneles del satélite. La tecla \textbf{1} controla el cohete. Estos son tres ejemplos de acciones iniciadas por teclado. Las naves rivales, criaturas, asteroides y el satélite se mueven continuamente. Fighter\_02 sale del hangar con la tecla \textbf{4}; la puerta se puede abrir sola con \textbf{5}. En el planeta rocoso, \textbf{6} reproduce el cohete por keyframes y \textbf{7} lo reinicia. El reloj visual permite apreciar el ciclo completo de día y noche, que dura 120 segundos; las luminarias se encienden juntas al oscurecer. El cielo cambia entre el espacio y las superficies.

\section{Problemas frecuentes}
\begin{description}[leftmargin=3cm,style=nextline]
\item[No abre el juego] Extraiga el ZIP completo. Ejecute \texttt{Verify-Release.ps1} y revise que la PC soporte OpenGL 3.3.
\item[No aparecen modelos o texturas] Abra \texttt{Jugar.cmd} desde la carpeta distribuida; restaure la estructura original del ZIP.
\item[No responden las teclas] Haga clic en la ventana. Use F9 para recuperar o liberar el cursor y compruebe el modo en pantalla.
\item[Una acción no se activa] Las acciones dependen de la ubicación. En R4 seleccione el cohete con Tab antes de pulsar 1; para abordar, acérquese a la nave.
\end{description}

\section*{Referencia de versión}
Este manual se preparó con el candidato \texttt{v0.4.0-rc.2}, commit \texttt{7ece353}, y capturas de pruebas del programa. Si cambia el ejecutable entregado, se deben revisar de nuevo los controles y las imágenes.
\end{document}
